\clearpage\begingroup\SourceLeaf{65}\HistoricalHead{60}{L}{PHYSIK. NACHLASS.}
\begin{center}[2.]\end{center}
\noindent{\small [S. 185]}
\begin{center}\textbf{Rechnung für die Wirkung eines Magnets \(SN\) in grosser Ferne \(AP=r\).}\end{center}
\begin{center}
\begin{tikzpicture}[x=.020cm,y=.020cm,line width=.55pt]
\coordinate (A) at (301,0);\coordinate (P) at (113,105);\coordinate (Q) at (48,219);
\draw (A)--(P);
\draw[dash pattern=on 2pt off 2.7pt] (P)--(-4,174);
\draw[dash pattern=on 2pt off 2.7pt,->] (P)--(Q);
\draw[dash pattern=on 2pt off 2.7pt] (A)--(235,111);
\draw[dash pattern=on 2pt off 2.7pt] (-10,0)--(267,0);
\draw (267,0)--(335,0);\draw (267,8)--(267,-7)--(335,-7)--(335,8);
\node[above,inner sep=1pt] at (Q) {$Q$};
\node[above right,inner sep=1pt] at (P) {$P$};
\node[above right,inner sep=1pt] at (235,111) {$P'$};
\node[below,inner sep=1pt] at (301,-7) {$A$};
\node[below,inner sep=1pt] at (267,-7) {$\scriptstyle N$};
\node[below right,inner sep=1pt] at (335,-7) {$\scriptstyle S$};
\node at (234,14) {$\scriptstyle u$};\node at (72,145) {$\scriptstyle v$};
\end{tikzpicture}
\end{center}
\noindent\begin{tabular}{@{}p{.10\textwidth}p{.86\textwidth}@{}}
\(u\)&Winkel zwischen \(AN\) und \(AP\)\\
\(PQ\)&Richtung der Kraft, welche in \(P\) auf positives magn. Fluidum wirkt\\
\(v\)&Winkel zwischen der Fortsetzung von \(AP\) und \(PQ\)
\end{tabular}\par
\noindent \(u\) und \(v\) werden nach einerlei Seite (beide von der linken nach der rechten
\begin{center}oder beide umgekehrt) gezählt\end{center}
\noindent\begin{tabular}{@{}p{.10\textwidth}p{.86\textwidth}@{}}
\(m\)&Moment des Magnetismus des Magnets \(SN\)\\
\(SN\)&seine magn[etische] Axe
\end{tabular}\par
\indent Dann ist
\[\operatorname{tang}v=\tfrac12\operatorname{tang}u\]
\noindent Stärke der Kraft
\[
\begin{aligned}
&\frac{m}{r^3}\surd(4\cos u^2+\sin u^2)\\
={}&\frac{m}{r^3}\surd(3\cos u^2+1)\\
={}&\frac{m}{r^3}\surd(4-3\sin u^2)\\
={}&\frac{m}{r^3}\surd(\tfrac52+\tfrac32\cos2u).
\end{aligned}
\]
\noindent Man kann diese Kraft auf drei verschiedene Arten in zwei zerlegen.
\par\medskip
\noindent\begin{tabular}{@{}p{.06\textwidth}cp{.66\textwidth}@{}}
I.&\(\dfrac{m}{r^3}\)&parallel mit \(NS\),\\[.6em]
&\(\dfrac{3m\cos u}{r^3}\)&nach Richtung \(AP\).
\end{tabular}\par\smallskip
\noindent{\small [S. 186]}\par\smallskip
\noindent\begin{tabular}{@{}p{.06\textwidth}ll@{}}
II.&\(\dfrac{2m}{r^3}\cos u\)&nach \(AP\),
\end{tabular}
\clearpage\SourceLeaf{66a}\HistoricalHead{61}{R}{BERECHNUNG MAGNETISCHER KRÄFTE.}
\noindent\begin{tabular}{@{}p{.06\textwidth}cp{.69\textwidth}@{}}
&\(\dfrac{m}{r^3}\sin u\)&nach der darauf senkrechten Richtung oder der Richtung, die
mit \(AN\) den Winkel \(u+90^0\) macht.\\[.6em]
III.&\(\dfrac{3m}{2r^3}\)&nach \(AP'\), wo \(NAP=PAP'\),\\[.6em]
&\(\dfrac{m}{2r^3}\)&parallel mit \(SN\).
\end{tabular}
\bigskip\begin{center}\textbf{BEMERKUNGEN.}\end{center}
{\small\noindent\hspace*{\parindent}Die hier mitgeteilten Formeln von \textsc{Gauss} sind zuerst von \textsc{Weber} in der Abhandlung \textit{Bemerkungen
über die Einrichtung und den Gebrauch des Bifilar-Magnetometers} in den Resultaten aus den Beobachtungen
des Magnetischen Vereins im Jahre 1837, S. 22 ff. mitgeteilt worden; dort findet sich auch die
erste der beiden Zeichnungen. \textsc{Gauss} selbst ist darauf zurückgekommen im Artikel 7 seiner \textit{Vorschriften
zur Bestimmung der magnetischen Wirkung, welche ein Magnetstab in der Ferne ausübt}, Resultate des
Magnetischen Vereins im Jahre 1840, S. 33; Werke V, S. 434.
\par\raggedleft\textsc{Schaefer}.\par}
\clearpage\endgroup
